\chapter{Конструкторская часть}

В данном разделе рассматривается устройство разрабатываемого программного
обеспечения. Описывается разбиение системы на основные модули, используемые
модели данных, а также порядок взаимодействия с выбранными большими
языковыми моделями при решении задачи проверки структуры PDF-документов.

\section{Общая структура программного обеспечения}

Проектируемое программное обеспечение состоит из следующих логических компонентов:
\begin{enumerate}
    \item модуль извлечения текстового содержимого из PDF-файла с сохранением
    координат строк (номер страницы и номер строки);
    \item модуль определения границ области <<СОДЕРЖАНИЕ/ОГЛАВЛЕНИЕ>>;
    \item модуль проверки выполнения условий варианта на основе регулярных выражений;
    \item модуль фиксации первого обнаруженного нарушения с указанием координат;
    \item модуль пакетной обработки PDF-файлов и формирования итогового отчёта
    в формате CSV.
\end{enumerate}

Такая модульная структура упрощает разработку, отладку и последующее расширение
программного обеспечения.

\section{Модель координат}

В качестве координат найденных совпадений используется пара значений:
номер страницы PDF-документа и номер строки на данной странице.
Оба значения нумеруются с единицы.

Использование данной модели координат позволяет однозначно указать место
обнаружения нарушения в документе без необходимости анализа геометрической
структуры PDF-файла, что соответствует требованиям задания.

\section{Использование регулярных выражений}

Для поиска и проверки структурных элементов документа применяются
регулярные выражения, учитывающие следующие особенности:
\begin{itemize}
    \item вариативность регистра символов;
    \item возможные пробелы между буквами в заголовках
    (например, <<В В Е Д Е Н И Е>>);
    \item наличие номерных префиксов
    (<<1>>, <<1.2>>), располагающихся перед заголовком;
    \item наличие ошибочных или недопустимых наименований
    (синонимы и близкие формулировки).
\end{itemize}

Данный подход позволяет корректно обрабатывать документы, оформленные
с отклонениями от шаблона, но сохраняющие требуемую структуру.

\section{Логика проверки варианта}

Алгоритм проверки одного PDF-файла включает следующие этапы:
\begin{enumerate}
    \item извлечение текстовых строк документа с сохранением их координат
    (номер страницы, номер строки);
    \item поиск границ области оглавления по заголовку
    <<СОДЕРЖАНИЕ/ОГЛАВЛЕНИЕ>> и характерной форме строк-элементов;
    \item проверку наличия в оглавлении строк
    <<ВВЕДЕНИЕ>>, <<ЗАКЛЮЧЕНИЕ>>,
    <<СПИСОК ИСПОЛЬЗОВАННЫХ ИСТОЧНИКОВ>>;
    \item проверку отсутствия корректных заголовков указанных разделов
    в основном тексте документа;
    \item проверку отсутствия ошибочных наименований разделов
    (например, <<Вступление>>, <<Выводы>>, <<Литература>>, <<Библиография>>)
    во всём документе;
    \item при обнаружении первого нарушения — фиксацию координат
    проблемной строки и завершение проверки файла.
\end{enumerate}

\section{Организация взаимодействия с БЯМ}

В рамках выполнения работы большие языковые модели использовались
для получения альтернативных реализаций алгоритма проверки структуры
PDF-документов. Взаимодействие с каждой моделью осуществлялось
в виде текстовых запросов, содержащих описание условий варианта
и требований к логике проверки.

Для всех моделей использовался одинаковый набор требований,
что позволило корректно сравнить полученные решения.
Результаты работы каждой модели представлены в виде программных
реализаций, приведённых в приложении~А.

\section{Анализ разработанных решений}

Реализация, полученная с использованием ChatGPT (GPT-5.2),
приведена в листинге~\ref{lst:chatgpt_core}.
Данная модель предложила наиболее полную и структурированную реализацию,
включающую извлечение текста из PDF-файла, формирование набора регулярных
выражений и базовые функции обработки строк.
Основное достоинство данного решения заключается в читаемости кода
и чётком разделении этапов обработки документа.

Реализация, полученная с использованием модели DeepSeek DeepThink,
представлена в листинге~\ref{lst:deepseek_core}.
В этом решении основной акцент сделан на алгоритм определения границ
области оглавления и проверку структуры документа.
Подход отличается более строгой логикой проверки и повышенной устойчивостью
к нетипичному форматированию PDF-файлов.

Реализация, полученная с использованием модели Qwen Coder,
приведена в листинге~\ref{lst:qwen_core}.
Данное решение отличается использованием оптимизации раннего выхода
из алгоритма при обнаружении всех обязательных элементов структуры.
Это позволяет сократить количество проверок и уменьшить время обработки
документа в ряде случаев.

\section{Листинги ключевых алгоритмов}

Полные реализации алгоритмов имеют значительный объём, поэтому в приложении
приведены укороченные листинги, содержащие ядро логики проверки структуры
документа без вспомогательного ввода-вывода и интерфейсных компонентов.
\begin{itemize}
    \item Листинг~\ref{lst:chatgpt_core} — реализация, полученная с использованием ChatGPT / GPT-5.2;
    \item Листинг~\ref{lst:deepseek_core} — реализация, полученная с использованием DeepSeek DeepThink;
    \item Листинг~\ref{lst:qwen_core} — реализация, полученная с использованием Qwen Coder.
\end{itemize}

\section{Вывод}

В данном разделе была рассмотрена архитектура программного обеспечения,
описана логика проверки структуры PDF-документов, а также проанализированы
решения, полученные с использованием различных больших языковых моделей.
